\documentclass[11pt,a4paper]{article}
\usepackage[margin=19mm,headheight=15pt]{geometry}
\usepackage{fontspec}
\setmainfont{Noto Serif}
\usepackage{amsmath,amssymb,mathtools}
\usepackage{xcolor}
\usepackage{fancyhdr}
\usepackage{enumitem}
\usepackage{microtype}
\setlength{\parindent}{0pt}
\setlength{\parskip}{6pt}
\definecolor{navy}{HTML}{183153}
\definecolor{lightgray}{HTML}{F3F5F7}
\pagestyle{fancy}\fancyhf{}
\lhead{\small ME 5309 -- Mid-Semester Examination}
\rhead{\small ME 5309}
\cfoot{\thepage}
\begin{document}
\begin{center}
{\Large\bfseries ME 5309 -- Mid-Semester Examination}\\[3pt]
{\large Answer Booklet}\\[8pt]
\end{center}
\textbf{Instructions:} Attempt any three questions. Each question carries 10 marks. Show working where appropriate.
\section*{Question 2 \hfill [10 marks]}
\subsection*{Q2(a) \hfill [2 marks]}

ẋ₁ = x₂ and ẋ₂ = (τ-b x₂+mgℓ sin θ)/J.
At equilibrium:\[\begin{aligned}τ_0 = -mg\ell sin \theta_e,\\ \text{an equilibrium exists for |τ₀| less than or equal to 2.4525 N m}\end{aligned}\]

\subsection*{Q2(b) \hfill [2 marks]}

\[\begin{aligned}A = [[0,1], [-mg\ell sin(60^{\circ})/J, -b/J]] = [[0,1], [-106.2, -2.5]],\\ B = [[0],[50]]\end{aligned}\]

\subsection*{Q2(c) \hfill [4 marks]}

At steady state the derivative terms vanish.
The proportional torque is Kp(θd-θ), which must balance mgℓ sin θ.

\subsection*{Q2(d) \hfill [2 marks]}

With gravity compensation, the gravity term cancels the plant's mgℓ sin θ term.
At rest, Kp(θd-θss)=0, so θss=θd and the steady-state error is zero.
The linearised gravity stiffness term is also cancelled, so the closed-loop poles generally change; the stiffness contribution mgℓ cos θd is removed.\par\smallskip The compensation term cancels the nominal gravity torque in the plant model. At steady state the proportional term must then be zero, which requires $\theta_d-\theta_{ss}=0$. This does not mean the transient response is unchanged: removing the gravity derivative changes the linearized stiffness.

\section*{Question 3 \hfill [10 marks]}
\subsection*{Q3(a) \hfill [2 marks]}

The estimate is noisy and can fluctuate from sample to sample.

\subsection*{Q3(b) \hfill [2 marks]}

\[\begin{aligned}y[k]=0.8818y[k-1]+0.0591x_k+0.0591x[k-1], DC gain 1\end{aligned}\]

\subsection*{Q3(c) \hfill [4 marks]}

With L=[ℓ₁,ℓ₂]ᵀ and C=[1,0], A-LC=[[-ℓ₁,1], [-ℓ₂,-2]].
Its characteristic polynomial is det(sI-(A-LC)) = (s+ℓ₁)(s+2)+ℓ₂ = s²+(ℓ₁+2)s+(2ℓ₁+ℓ₂).
Desired poles at -10 give (s+10)²=s²+20s+100.
Hence ℓ₁+2=20 gives ℓ₁=18, and 2ℓ₁+ℓ₂=100 gives ℓ₂=64.\par\smallskip Expanding $(s+\ell_1)(s+2)+\ell_2$ gives $s^2+(\ell_1+2)s+(2\ell_1+\ell_2)$. The desired polynomial is $(s+10)^2=s^2+20s+100$. Matching the $s$ coefficient gives $\ell_1=18$; matching the constant term then gives $36+\ell_2=100$, hence $\ell_2=64$.

\subsection*{Q3(d) \hfill [2 marks]}

The observer is better and the filter is safer in some cases.

\section*{Question 4 \hfill [10 marks]}
\subsection*{Q4(a) \hfill [2 marks]}

Poles are -2±2j and Nff=8.

\subsection*{Q4(b) \hfill [2 marks]}

The actual feedforward contribution is 0.75x8=6.

\subsection*{Q4(c) \hfill [4 marks]}

The desired polynomial is (s+4)(s²+4s+8).
The closed-loop polynomial should be cubic because an integrator is added.

\subsection*{Q4(d) \hfill [2 marks]}

At steady state, a stable closed loop with integral action drives ξ̇=r-θ to zero; therefore r-θ=0 and the tracking error is zero, assuming no saturation or other implementation limits prevent convergence.
A practical failure mode is integrator windup during actuator saturation; the standard remedy is anti-windup, such as back-calculation or conditional integration.\par\smallskip With integral action, a constant tracking error keeps accumulating in the integral state, so the controller can adjust its output until the error goes to zero, assuming the closed loop is stable and the actuator is not constrained. Under saturation the integral state can keep growing; anti-windup limits or back-calculates that state.

\end{document}